\setcounter{section}{1}
\setcounter{theorem}{5}
\section{Preliminaries}
\begin{definition}[Distance measure]
For any two vectors $a,b\in\R^d$, the distance between them in the $\ell_2$-metric is denoted by $\|a-b\|=\left(\sum_{i=1}^d|a_i-b_i|^2\right)^{1/2}$.
Their distance in the $\ell_0$-metric (Hamming distance) is denoted by 
$\|a-b\|_0=|\{i\in[d]:a_i\neq b_i\}|$, i.e.,
the number of coordinates on which $a$ and $b$ differ.
\end{definition}


\subsection{Quantum query model}

We consider the quantum query model in this work. Let $X:=\{x_1,\dots,x_n\}$ be a set of $n$ input points and $\ora_X$ be the corresponding oracle. We can access the $i$-th data point $x_i$ by making the query 
\begin{align}
    \ket{i}\ket{0} \xrightarrow{\ora_X} \ket{i}\ket{x_i},  
\end{align}
and we can make queries to elements in $X$ in superposition. Note that $\ora_X$ is an unitary transformation in the formula above. Hence, a quantum algorithm with access to $\ora_X$ can be represented as a sequence of unitary transformations. 

Consider a quantum algorithm $\A$ with access to an oracle $\ora$ and a initial state $\ket{0}:= \ket{0}_{Q}\ket{0}_{A}\ket{0}_W$, where the registers $Q$ and $A$ are for the queries and the answers from the oracle, and the register $W$ is the working space which is always hold by $\A$.  Then, we can represent the algorithm as 
\begin{align}
    U_T \ora U_{T-1}\cdots\ora U_1  \ket{0}. 
\end{align}
Let $\ket{\psi}_i= U_i\ora \cdots \ora U_1\ket{0}:= \sum_{i,z} \ket{i}_Q\ket{0}_A\ket{z}_W$ be the state right before applying the $i$-th $\ora$, then
\begin{align}
    \ora \ket{\psi}_i:= \sum_{i,z} \ket{i}_Q\ket{x_i}_A\ket{z}_W.
\end{align}
\subsection{Quantum subroutine for unstructured searching and minimum finding}

\begin{definition}[Unstructured search]
Given a set $P$ of $n$ elements in $\{0,1\}$, decide whether there exists a $1$ in $P$. 
\end{definition}

\begin{theorem}[Grover's search algorithm~\cite{r24,r37}]\label{thm:grover_algo}
There is a quantum algorithm for unstructured search with running time $O(\sqrt{n})$.
\end{theorem}

By \cref{thm:grover_algo} and BBBV's argument \cite{r9}, the quantum time complexity of unstructured search is $\Theta(\sqrt{n})$. We can also get a $\tilde{O}(\sqrt{n})$ quantum algorithm for minimum finding by combining Grover's search algorithm and binary search.
\begin{theorem}[Quantum minimum finding~\cite{r23}]\label{thm:mim}
There is a quantum algorithm that finds from a set of $n$ elements with values in $\mathbb{R}$, the index of the minimum element of the set, with success probability $\frac{1}{2}$ and run time $\tilde{O}(\sqrt{n})$.
\end{theorem}

\subsection{Problem definitions}
In this subsection, we first formally define  $\OV$, $\CP$, and $\BCP$. Then we show the folklore algorithms for $\CP$, $\BCP$, and $\OV$ by Grover's algorithm, which run in time $\tilde{O}(n)$. 
\begin{definition}[Orthogonal Vectors, $\OV$]\label{def:OV}
Given two sets $A,B$ of $n$ vectors in $\mathbb \{0,1\}^d$ as input, find a pair of vectors $a\in A$, $b\in B$ such that
$\ip{a}{b}=0$, where the inner product is taken in $\mathbb{Z}$.\footnote{Our definition is slightly different than some of the literature, for example, \cite{r18}, which is searching among pairs inside one set. Those two definitions are equivalent up to constant in complexities.}
\end{definition}

We denote $\OV$ with input length $n$ and dimension $d$ as $\OV_{n,d}$. We will use this notation when we need to specify the parameters in the following sections.  

\begin{definition}[Closest Pair Problem, $\CP$]
Given a set $P$ of $n$ points in $\mathbb R^d$ and a distance measure $\Delta$, find a pair of distinct points $a,b\in P$ such that $\Delta(a,b)$ is the smallest among all distinct pairs in $P$.
\end{definition}
Similar to $\OV$, we denote $\CP$ with input length $n$ and dimension $d$ as $\CP_{n,d}$. We will use this notation when the parameters in the following sections are required to be specified. Note that in this work, we consider $\Delta(a, b) = \norm{a-b}$ as the distance measure for \textsf{CP} and \textsf{BCP}.

\begin{definition}[Bichromatic Closest Pair Problem, \textsf{BCP}]
Given two sets $A,B$ of $n$ points in $\mathbb R^d$ and a distance measure $\Delta$, find a pair of points $a\in A$, $b\in B$ such that 
\begin{align}
    \Delta(a,b)= \min_{a\in A,b\in B} \Delta(a,b).
\end{align}
\end{definition}
We also define an approximate version of $\BCP$ as follows.
\begin{definition}[$(1+\xi$)-approximate Bichromatic Closest Pair Problem, $\xiBCP$]
Given two sets $A,B$ of $n$ points $\in \mathbb R^d$ and a distance measure $\Delta$, find a pair of points $a\in A$, $b\in B$ such that 
\begin{align}
    \Delta(a,b) \leq (1+\xi)\min_{a\in A,b\in B} \Delta(a,b).
\end{align}
\end{definition}
Same as $\CP$, we use $\BCP_{n,d}$ and $\xiBCP_{n, d}$ to specify the parameters. 

\begin{definition}[Element Distinctness Problem, \textsf{ED}]
  Let $f:[n] \rightarrow [m]$ be a given function. Decide whether there exist distinct $i, j\in[n]$ such that $f(i) = f(j)$.
\end{definition}

For this problem, Ambainis \cite{r5} gave a quantum algorithm with time complexity $\tilde{O}(n^{2/3})$, which matches the lower bound proved by Aaronson and Shi \cite{r1} up to a polylogarithmic factor.

\begin{theorem}\label{cor:trivial_alg}
There are $\tilde{O}(n)$-time quantum algorithms for $\CP$ and $\BCP$ when $d=O(\poly\log n)$.
\end{theorem}
\begin{proof}
We can solve $\CP$ and $\BCP$ by searching the minimum distance through all pairs by the algorithm of \cref{thm:mim}. There are $O(n^2)$ pairs and checking each pair took $O(d)$ time, so the total running time is $O(nd)$. For $d=O(\poly\log n)$, the time complexity equals to $\tilde{O}(n)$.
\end{proof}

\subsection{Fine-grained complexity}

As we have mentioned earlier in the introduction, a fine-grained reduction from problem $\mathsf{P}$ to $\mathsf{Q}$ with conjectured lower bounds $p(n)$ and $q(n)$, respectively, has the property that if we can improve the $q(n)$ time for $\mathsf{Q}$, then we can also improve the $p(n)$ time for $\mathsf{P}$. We give the formal definition by Williams~\cite{r46} in below. 

\begin{definition}[Fine-grained reduction, \cite{r46}]\label{def:fg}
Let $p(n)$ and $q(n)$ be non-decreasing functions of $n$. Problem $\mathsf{P}$ is ($p,q$)-reducible to problem $\mathsf{Q}$, denoted as $(\mathsf{P},p)\leq_{\FG} (\mathsf{Q},q)$, if for every $\epsilon$, there exist $\delta>0$, an algorithm $R$ for solving $\mathsf{P}$ with access to an oracle for $\mathsf{Q}$, a constant $d$, and an integer $k(n)$, such that for every $n\ge 1$, the algorithm $R$ takes any instance of $\mathsf{P}$ of size $n$ and 
\begin{itemize}
    \item $R$ runs in at most $d\cdot (p(n))^{1-\delta}$-time, 
    \item $R$ produces at most $k(n)$ instances of $\mathsf{Q}$ adaptively, that is, the $j$th instance $X_j$ is a function of $\{(X_i,y_i)\}_{1\leq i<j}$ where $X_i$ is the $i$th instance produced and $y_i$ is the answer of the oracle for $\mathsf{Q}$ on instance $X_i$, and 
    \item the sizes $n_i$ of the instances $X_i$ for any choice of oracle answers $y_i$ obeys the inequality
    \begin{align}\label{eq:fgr}
        \sum_{i=1}^{k(n)} (q(n_i))^{1-\epsilon} \leq d\cdot (p(n))^{1-\delta}.
    \end{align}
\end{itemize}
\end{definition}

Let $(\mathsf{P},p)\leq_{\FG} (\mathsf{Q},q)$ for some non-decreasing function $p(n)$ and $q(n)$. If for every $\eps > 0$, we can solve problem $\mathsf{Q}$ in time $q(n)^{1-\epsilon}$ with probability $1$ for all input length $n$, then there exists a $\delta>0$ such that we can solve the problem $\mathsf{P}$ in time $p(n)^{1-\delta}$ 
by \cref{eq:fgr}.

Here are some known results about fine-grained reductions.
\begin{theorem}[\cite{r29,r49}]\label{thm:classical_fg}
\begin{align}
    (\mathsf{CNF\text{-}SAT}_{n},2^n)\leq_{\FG} (\OV_{n_1,d_1},n_1^2)\leq_{\FG} (\BCP_{n_2,d_2},n_2^2)\leq_{\FG} (\CP_{n_3,d_3},n_3^2),
\end{align}
where $d_1=\Theta(\log n_1)$, $d_2=\Theta(\log n_2)$ and $d_3=(\log  n_3)^{\Omega(1)}$.
\end{theorem}

\begin{remark}
The second reduction from $\OV$ to $\BCP$ has been improved to $d_2 = 2^{O(\log^* n)}$ by Chen~\cite{r17}.
\end{remark}

There are several plausible hypotheses in fine-grained complexity, which can imply conditional hardness results for many interesting problems. We first give the definition of the strong exponential time hypothesis (SETH). 
\begin{hypothesis}[Strong Exponential Time Hypothesis, \textsf{SETH}]
For every $\epsilon > 0$, there exists a $k = k(\epsilon) \in \N$ such that no algorithm can solve $k$-$\SAT$ (i.e., satisfiability on a CNF of width $k$) in $O(2^{(1 - \epsilon)m})$ time where $m$ is the number of variables. Moreover, this holds even when the number of clauses is at most $c(\epsilon) \cdot m$ where $c(\epsilon)$ denotes a constant that depends only on $\epsilon$.
\end{hypothesis}

Another popular conjecture is the orthogonal vector hypothesis (\textsf{OVH}):
\begin{definition}[Orthogonal Vector Hypothesis, \textsf{OVH}]
For every $\epsilon > 0$,  
there exists a $c\ge 1$ such that $\OV_{n,d}$ requires $n^{2-\epsilon}$ time when $d=c\log n$.
\end{definition}

\begin{remark}
Under SETH, we can have the following conclusions from \cref{thm:classical_fg}: 
\begin{itemize}
    \item $\OVH$ is true. 
    \item For all $\epsilon>0$, there exists a $c>0$ such that $\BCP_{n,c\log n}$ cannot be solved by any randomized algorithm in time $O(n^{2-\epsilon})$.
    \item For all $\epsilon>0$, there exists a $c>0$ such that $\CP_{n,(\log n)^c}$ cannot be solved by any randomized algorithm in time $O(n^{2-\epsilon})$.
\end{itemize}
\end{remark}

\subsection{The framework for quantum walk search}
\label{sec:qw-framework}
In this subsection, we review the quantum walk framework for the Markov chain search problem and demonstrate how to use it to solve the element distinctness problem. For simplicity, we use the transition matrix $P$ to refer to a Markov chain, where $P = (p_{xy})_{x, y \in X}$ for $X$ being the state space of $P$ and $p_{xy}$ being the transition probability from $x$ to $y$. An irreducible and ergodic Markov chain has a unique stationary distribution $\pi$, which is also the unique eigenvector of $P$ with eigenvalue 1. Let $M \subseteq X$ be a set of marked elements. In the Markov chain search problem, the objective is to find an $x \in M$. We can perform the following actions: setup, sampling from the $\pi$ with cost $\mathsf{S}$; update, making a transition with cost $\mathsf{U}$, and checking whether the current state is marked or not with cost $\mathsf{C}$. To solve the search problem classically, we perform a random walk as follows. Start from a point sampled from $\pi$ and check if it is marked. If not, make a number of transitions on $P$ until it mixes, and then check again. We then repeat this process until a marked state is found. The cost of this random walk algorithm is $O(\mathsf{S} + \frac{1}{\varepsilon}(\frac{1}{\delta}\mathsf{U} + \mathsf{C}))$, where $\varepsilon:= |M|/|X|$ and $\delta$ is the spectral gap of $P$.

Quantum analogues of random walks, namely, quantum walks, have been developed for solving different problems. In 2003, Ambainis~\cite{r5} proposed a quantum walk algorithm for solving the element distinctness problem. His algorithm also solves the Markov chain search problem on the Johnson graph with cost $O(\mathsf{S} + \frac{1}{\sqrt{\varepsilon}}(\frac{1}{\sqrt{\delta}}\mathsf{U} + \mathsf{C}))$. In 2004, Szegedy~\cite{r45} gave a quantum walk algorithm for more generalized Markov chains with cost $O(\mathsf{S} + \frac{1}{\sqrt{\varepsilon\delta}}(\mathsf{U} + \mathsf{C}))$. We can view Szegedy's quantum walk as a quantum counterpart of a random walk, where one checks the state after each transition. Szegedy's quantum walk only detects the presence of a marked state, but cannot find one without extra costs. In 2006, Magniez et al.~\cite{r33} proposed a quantum walk search framework that unified the advantages of the quantum walks in~\cite{r5} and~\cite{r45}. In this quantum walk framework, we can perform the following operations:
\begin{itemize}
  \item \textbf{Setup:} with cost $\mathsf{S}$. preparing the initial state $\ket{\pi} = \frac{1}{\sqrt{|X|}}\sum_{x}\sqrt{\pi_x}\ket{x}$.
  \item \textbf{Update:}  with cost $\mathsf{U}$. applying the transformation $\ket{x}\ket{0} \mapsto \ket{x}\sum_{y\in X}\sqrt{p_{xy}}\ket{y}$.
  \item \textbf{Checking:} with cost $\mathsf{C}$, applying the transformation: $\ket{x} \mapsto \big\{\begin{smallmatrix} -\ket{x} \quad &\text{if $x\in M$} \\ \ket{x} \quad &\text{otherwise.}\end{smallmatrix}$
\end{itemize}
The main result of \cite{r33} is summarized as follows.
\begin{lemma}[\cite{r33}]
  \label{lem:framework}
  Let $P$ be an irreducible and ergodic Markov chain $P$ on $X$. Let $M \subseteq X$ be a subset of marked elements. Let $\varepsilon: = |M|/|X|$ and $\delta$ be the spectral gap of $P$. Then, there exists a quantum algorithm that with high probability, determines $M$ is empty or finds an $x \in M$ with cost $O(\mathsf{S} + \frac{1}{\sqrt{\varepsilon}}(\frac{1}{\sqrt{\delta}}\mathsf{U} + \mathsf{C}))$.
\end{lemma}

To solve the element distinctness problem, we define a Markov chain, following the work~\cite{r5,r11,r28}. The state space $X$ is all subsets of $[n]$ with size $r$. The Markov chain is based on the Johnson graph on $X$, where an edge is connecting $S$ and $S'$ if and only if $|S \cap S'| = r-1$. The transition probability on each edge is hence $\frac{1}{r(n-r)}$. A state $S$ is marked when there exist distinct $i, j \in S$ such the $i^{th}$ and the $j^{th}$ items are the same. The Markov chain has spectral gap $\delta \geq 1/r$ (see~\cite{r28}) and it is easy to verify that $\varepsilon \geq \binom{n-2}{r-2}/\binom{n}{r} = O(r^2/n^2)$. If we only consider the query complexity, the setup procedure costs $r$ queries, the update procedure costs one query, and the checking procedure does not cost any query. Choosing $r = n^{2/3}$ yields the optimal query complexity $O(n^{2/3})$.


